%% ma: L10-B02-0037
%% lop: 10
%% chuong: 1
%% bai: 2
%% dang: 10.2.4
%% loai: DS
%% muc: [NB, NB, NB, NB]
%% dap_an: "ĐSĐĐ"
%% nguon: "Ảnh giáo viên gửi 10/10/2026 (ThS. Nguyễn Hoàng Việt, Chương 2 Tập hợp), Đề 2 (đề thứ hai, không ghi tiêu đề), Phần 2 Câu 15"
%% kien_thuc: "Bài 2 (tập hợp và các phép toán trên tập hợp)"
%% trang_thai: cho-duyet
%% ly_do: "Dạng toán mới đề xuất 10.2.4: tập con và số tập con"
\begin{ex}
Giả sử $A=\{2;4;6\}$, $B=\{2;6\}$, $C=\{4;6\}$, $D=\{4;6;8\}$. Vậy:
\choiceTF
{\True $B\subset A$.}
{$A\subset B$.}
{\True $C\subset A$.}
{\True $C\subset D$.}
\loigiai{a) Mọi phần tử của $B$ ($2,6$) đều thuộc $A$: đúng. b) $4\in A$ nhưng $4\notin B$: sai. c) $4,6\in A$: đúng. d) $4,6\in D$: đúng.}
\end{ex}
